\begin{textsample}{POS Dim 1 – ac – Score 40.00 – ac/ac\_em\_18.txt}  \label{ex:f1_pos_192}
11.1.
Sistemáticas de acompanhamento e avaliação Na era da tecnologia, da globalização e dos avanços científicos em \textbf{que} \textbf{vivemos}, a Educação é \textbf{central} no processo de inserção e interação social dos \textbf{sujeitos}, portanto, no \textbf{que} \textbf{diz} respeito à sua qualidade de oferta, os processos avaliativos são de suma importante, pois \textbf{configura} -se como \textbf{um} dos pilares fundamentais do trabalho pedagógico da escola e da sala de aula.
Avaliação e aprendizagem devem caminhar juntas no processo educativo.
Dessa forma, não há avaliação sem aprendizagem, assim como não há aprendizagem sem avaliação.
Nessa mesma \textbf{direção}, Santos e Gontijo afirmam \textbf{que} : > [... ] avaliar é \textbf{um} processo natural \textbf{que} faz parte do repertório das ações dos seres humanos e está indissociavelmente ligado à busca de todo tipo de ação \textbf{que} provoque mudanças, mudanças essas \textbf{que} indicam respostas concretas ao desejo por “ algo melhor ”.
O ato de avaliar tem como foco a construção dos melhores resultados possíveis.
O ser humano permanentemente avalia em busca de \textbf{uma} melhor qualidade de vida, da construção de si mesmo e do melhor modo de ser e \textbf{viver} ( \textbf{SANTOS} ; GONTIJO, 2018, p. 31 ).
Portanto, a avaliação não deve ser somente o momento da realização das \textbf{provas} e testes, mas \textbf{um} processo contínuo e \textbf{que} ocorre dia após dia, visando a identificação dos objetivos de ensino alcançados, percebendo as dificuldades do aluno em seus movimentos para a constituição de conhecimentos e o domínio de habilidades, previstas para a sua etapa escolar.
Nesse sentido, a forma avaliativa funciona como \textbf{um} elemento de integração e motivação para o processo de \textbf{ensino-aprendizagem}.
A avaliação apresenta -se em várias modalidades, entre as quais está a avaliação somativa ou classificatória \textbf{que}, conforme Haydt ( 2000 ), tem como função classificar os alunos ao final da unidade, semestre ou ano letivo, segundo níveis de aproveitamento apresentados.
O objetivo da avaliação somativa é classificar o aluno para determinar se ele será aprovado ou reprovado e está vinculado à noção de medir.
O sistema educacional, muitas vezes, tem se apoiado na avaliação classificatória com a pretensão de verificar aprendizagem ou competências através de medidas, de quantificações.
Este tipo de avaliação pressupõe \textbf{que} as pessoas aprendam do mesmo modo, nos mesmos momentos e tentam evidenciar competências isoladas.
Ou seja, algumas pessoas, \textbf{que} por diversas razões têm maiores condições de aprender, aprendem mais e melhor.
Outras, com características diferentes, \textbf{que} não respondem tão bem ao conjunto de \textbf{disciplinas}, aprendem cada vez menos e são muitas vezes excluídas do processo de escolarização.
Dentre as modalidades de avaliação, há também outras duas \textbf{que} são : a avaliação formativa e a diagnóstica.
A avaliação formativa se constitui em \textbf{um} processo complexo, cujo detalhamento se dará na relação direta entre professores e estudantes, ou seja, deve ser dimensionada e modulada para cada realidade escolar, não havendo \textbf{uma} fórmula a ser aplicada, mas o desenvolvimento de \textbf{um} processo \textbf{que} é parte da própria aprendizagem do estudante e do fazer pedagógico do professor.
A referida modalidade de avaliação é chamada formativa no sentido de \textbf{que} indica como os alunos estão se mobilizando em \textbf{direção} aos objetivos visados.
Assim, a avaliação formativa tem como função informar o aluno e o professor sobre os resultados \textbf{que} estão sendo alcançados durante o desenvolvimento das atividades, subsidiando -se em melhorar o ensino e a aprendizagem ; localizar, apontar, discriminar deficiências, insuficiências, no desenvolvimento do ensino-aprendizagem para eliminá -las ; proporcionar feedback de ação ( leitura, explicações, exercícios ) ( SANT’ANNA, 2001, p. 34 ).
A avaliação formativa adota vários tipos de \textbf{métodos} avaliativos, sempre de maneira engajada, oferecendo ao aluno a oportunidade de ser coautor da construção do próprio conhecimento e, portanto, utiliza algumas ferramentas, como as seguintes : - Autoavaliação - Prerrogativa para estimular as competências socioemocionais, motivando a empatia e a sinceridade, estimula o protagonismo dos alunos e a prática da autogestão. - Testes tradicionais - O \textbf{erro} é \textbf{visto} como parte do processo e não como falta grave. - Simulados - Preparam os estudantes para determinados exames, estimulando \textbf{um} estudo com mais perseverança. - Seminários - Promove autonomia, protagonismo, estimula a participação ativa. - Trabalhos em grupo - Os alunos aprendem a lidar com a opinião do outro e com a diversidade e incentiva a empatia.
Vale \textbf{lembrar} \textbf{que} esse tipo de avaliação visa formar e informar os alunos a respeito de seu desenvolvimento como cidadãos cientes de seu lugar no mundo.
Portanto, os \textbf{métodos} avaliativos \textbf{que} constituem \textbf{uma} avaliação formativa podem ser abertos ao \textbf{método} didático do professor, mas é \textbf{imprescindível} \textbf{que} a autoavaliação esteja presente sempre.
Portanto, não há melhor ou pior avaliação ; cada avaliação possui objetivos distintos.
O \textbf{ideal}, no entanto, é usá -las em diferentes circunstâncias, unindo o melhor de cada modelo avaliativo.
As avaliações somativas são essenciais para informar e situar os estudantes da escola como \textbf{um} todo.
Todavia, a avaliação formativa é utilizada para \textbf{um} processo contínuo e longo em \textbf{que} o \textbf{erro} não é nada mais \textbf{que} \textbf{um} fator \textbf{inerente} ao processo de \textbf{ensino-aprendizagem} dos estudantes.
Outro perfil avaliativo é a avaliação diagnóstica, constituída por \textbf{uma} sondagem, projeção e retrospecção da situação de desenvolvimento do aluno, dando -lhe elementos para verificar o \textbf{que} aprendeu e como aprendeu.
É \textbf{uma} etapa do processo educacional \textbf{que} tem por objetivo verificar em \textbf{que} medida os conhecimentos anteriores ocorreram e o \textbf{que} se faz necessário planejar para selecionar dificuldades encontradas.
A avaliação diagnóstica contribui para a compreensão das particularidades de aprendizagem de cada indivíduo e orienta sobre o \textbf{direcionamento} mais assertivo sobre o \textbf{que} deverá ser feito a partir do \textbf{que} foi diagnosticado ; possui ação preventiva, promove o conhecimento de aptidões, interesses e capacidades para direcionar planejamentos futuros, \textbf{revelando} causas e dificuldades de aprendizagens.
Assim, seus resultados servem para identificar e propor novos procedimentos de \textbf{ensino-aprendizagem}.
Conforme a realidade escolar do docente, este tem a liberdade para realizar as diferentes modalidades de avaliação relatadas, fazendo isso no momento do planejamento de suas aulas.
A avaliação formativa possui \textbf{um} papel \textbf{central} no \textbf{que} se refere à organização do currículo para o ensino médio, \textbf{embasado} na repartição do tempo escolar em duas partes indissociáveis - Formação Geral Básica e Itinerários Formativos, ambas articuladas pelo Projeto de Vida.
Nesse sentido, a avaliação formativa atua como instrumento \textbf{metodológico} essencial para \textbf{que} sejam atingidos os objetivos de aprendizagem propostos no currículo.
Isso visa manter os pressupostos \textbf{que} levaram à elaboração deste documento e adequá -los, naquilo \textbf{que} for necessário, à nova legislação vigente. \textbf{Contudo}, além de conhecer diferentes tipos de avaliação, é necessário também levar em conta os \textbf{métodos} de ensino \textbf{que} precedem as avaliações.
Nesse sentido, as metodologias ativas — práticas pedagógicas pensadas em conformidade com o \textbf{que} prega a BNCC — vem para modificar \textbf{um} pouco o ensino tradicional e dar base para \textbf{que} os estudantes tenham autonomia para trilhar o caminho rumo ao aprendizado, de maneira ativa.
As metodologias ativas, junto às avaliações formativa e somativa, são realmente \textbf{interessantes} e devem ser pensadas e desenvolvidas juntas, e não de maneira dissociada.
Nos últimos anos, o ensino médio tem passado por \textbf{inúmeras} mudanças, tendo seus marcos legais definidos por documentos \textbf{norteadores}, dentre os quais cabe destacar a Lei 13.415/2017 \textbf{que}, em seu Art. 3º, acrescentou o art. 35 -A, ampliando o texto da LDB 9394/96, e ressaltando ainda \textbf{que} : > § 7º Os currículos do ensino médio deverão considerar a formação integral do aluno, de maneira a adotar \textbf{um} trabalho voltado para a construção de seu projeto de vida e para sua formação nos aspectos físicos, cognitivos e socioemocionais. > § 8º Os conteúdos, as metodologias e as formas de avaliação processual e formativa serão organizados nas redes de ensino por meio de atividades \textbf{teóricas} e práticas, \textbf{provas} orais e escritas, seminários, projetos e atividades on-line, de tal forma \textbf{que} ao final do ensino médio o \textbf{educando} demonstre : > I - domínio dos princípios científicos e tecnológicos \textbf{que} presidem a produção \textbf{moderna} ; > II - conhecimento das formas \textbf{contemporâneas} de linguagem ( BRASIL, 2017, p. 2-3 ).
O Currículo de Referência Único do Acre propõe às escolas públicas da Rede de Educação Básica do Estado do Acre \textbf{um} processo avaliativo, orientado, normatizado e instrumentado pela instrução normativa nº 01 de 28 de fevereiro de 2019, publicada na edição nº 12.504 do D.O.E. de 01 de março de 2019, página 65, \textbf{que} considera o disposto na Lei de Diretrizes e Bases da Educação Nacional nº 9394/96, no Parecer do Conselho Nacional de Educação CNE/CEB nº 7/2010, na Resolução CNE/CEB nº 4/2010, na Resolução CNE/CEB nº 7/2010, no Parecer nº 15/2001 do Conselho Estadual de Educação – CEE e nas demais normas vigentes e \textbf{que} ainda serão elaboradas : > Art. 1{\EmojiFont °} Instruir o processo de avaliação da aprendizagem dos alunos do Ensino Fundamental e Médio das Escolas da Rede Pública de Educação do Estado do Acre \textbf{face} à interpretação dos instrumentos legais regulamentadores da matéria e assim estruturá -lo : > I. o ano letivo, para efeitos de organização e critérios \textbf{que} medem resultados, está dividido em 4 ( quatro ) bimestres ; > II. ao final de cada bimestre o aluno precisará ter alcançado \textbf{uma} nota resultante da aplicação de diversos instrumentos da avaliação formativa da aprendizagem tais como : \textbf{provas} e testes, trabalhos individuais ou em grupos, pesquisa, leituras complementares, diagnósticos, \textbf{tarefas} para a casa, exercícios desenvolvidos em sala de aula, tendo o alcance da média de qualidade \textbf{que} é 7,0 ( sete ).
A educação escolar orienta -se, intencionalmente, por metas \textbf{que} objetivam acompanhar todo o processo de ensino e aprendizagem.
Tais intenções da ação educativa só adquirem sentido quando levadas em consideração a natureza social e a função socializadora da educação escolar, as quais têm como razão primordial a promoção do desenvolvimento humano.
Incluem -se nesse processo os procedimentos didáticos assumidos pelo professor e a avaliação como ferramenta fundamental na aprendizagem ( DARSIE, 1996 ).
Como já comentado, avaliação e aprendizagem devem caminhar juntas.
Sendo assim, o professor, conforme a sua realidade escolar, poderá escolher as estratégias avaliativas, levando em consideração as propostas de atividades/metodológicas utilizadas para mobilizar o desenvolvimento de competências e habilidades, isto é, o professor poderá fazer uso de recursos avaliativos diferenciados.
Segundo Santos e Varela ( 2007 ), ao avaliar, o professor deve utilizar técnicas diversas e instrumentos variados, para \textbf{que} assim \textbf{consiga} diagnosticar o começo, o durante e o fim de todo o processo avaliativo percorrido pelo \textbf{discente}, para \textbf{que}, a partir de então, possa progredir no processo didático e retomar o \textbf{que} foi insatisfatório para o processo de aprendizagem dos educandos.
Por fim, os estudantes de ensino médio, para alcançarem determinado aprendizado, deparam -se com \textbf{um} processo complexo, de cunho pessoal, para o qual os professores e a escola podem contribuir ao permitir a comunicabilidade, ao situá -los em seu grupo e a oportunizar \textbf{que} aprendam a respeitar e a fazer -se respeitar.
Para tal, é salutar \textbf{que} professores e escolas criem situações em \textbf{que} o estudante deva sentir -se desafiado ou \textbf{instigado} a participar e questionado pelo jogo do conhecimento, adquirindo o \textbf{espírito} de pesquisa e o desenvolvimento de capacidades de \textbf{raciocínio} e de autonomia ( BRASIL, 1999 ).
Assim, os parâmetros de avaliação voltados para a etapa do ensino médio devem estar em consonância com os princípios \textbf{que} \textbf{permeiam} esta etapa de ensino, de modo a subsidiar o professor para diagnosticar e monitorar as aprendizagens, identificar as potencialidades e eventuais dificuldades presentes no processo, além de acompanhar os resultados das práticas de ensino.
Desse modo, avaliação deverá ter como referência : - O desenvolvimento integral do aluno, de modo a garantir a autonomia e o apropriação das competências socioemocionais. - Flexibilização de metodologias para efetivo desenvolvimento das competências cognitivas e socioemocionais. - Coerência entre a prática pedagógica e os processos avaliativos. - A avaliação dos resultados das avaliações da aprendizagem, de modo a proporcionar reflexões acerca das evidências e diagnósticos para os ajustes necessários nas práticas pedagógicas. - A avaliação formativa deve está constantemente presente como instrumento, de modo a garantir a aferição do desenvolvimento das competências gerais e específicas do Currículo. - Os instrumentos avaliativos devem ser diversificados e articulados entre si. - As avaliações devem acompanhar a flexibilização do ensino médio e apoiar os jovens em seus Projetos de Vida.

% matched lemmas: central, configurar, conseguir, contemporâneo, contudo, direcionamento, direção, discente, disciplina, dizer, educar, embasar, ensino-aprendizagem, erro, espírito, face, ideal, imprescindível, inerente, instigar, interessante, inúmero, lembrar, metodológico, moderno, método, norteador, permear, prova, que, raciocínio, revelar, santo, sujeito, tarefa, teórico, um, ver, vivar, viver
\end{textsample}
